\chapter{Post-Quantum Migration of Blockchains}
\label{ch:blockchains}
\chaptermark{Post-Quantum Migration of Blockchains}

\section{Why this chapter matters}
Blockchains are one of the hardest post-quantum migration problems in existence, for two reasons the previous chapter only touched. First, the funds are secured by \emph{signatures} whose public keys are, sooner or later, \emph{permanently public} on an immutable ledger -- so ``harvest now, decrypt later'' becomes a much sharper ``store now, steal later.'' Second, there is no administrator: changing the rules of Bitcoin or Ethereum is a governance problem settled by rough consensus and forks, not a software update pushed to servers. This chapter looks at how three ecosystems are actually approaching the transition, as of mid-2026: Bitcoin and Ethereum, the two largest, and Cosmos, which is furthest along in shipped code. Everything here is a moving target: the concrete proposals are drafts, not deployed consensus, and the reader should treat dates and statuses as a snapshot.

Mosca's inequality $x + y > z$ (Chapter~\ref{ch:quantum}) states the problem in
one line, and blockchains are the case where it degenerates. The shelf life $x$
of an exposed on-chain public key is not measured in years: the ledger is
immutable and public, so the key is a target for as long as the chain exists and
the coins remain unmoved. With $x$ unbounded, $x + y > z$ holds for \emph{any}
migration time $y$ and \emph{any} finite $z$. The inequality therefore cannot be
satisfied by better hardware forecasts or by moving faster; it can only be
escaped by removing the exposure --- which means moving the coins. That is why
every roadmap in this chapter, whatever else it disputes, ends up arguing about
how to make users act.

\section{Learning goals}
After reading this chapter, you should be able to:

\begin{itemize}
  \item explain which blockchain primitives a quantum computer breaks (Shor) versus merely weakens (Grover), and the public-key-exposure principle;
  \item describe Bitcoin's exposure by output type and the current draft proposals (BIP-360, BIP-361), and distinguish exposure caused by \emph{protocol design} from exposure caused by \emph{user behaviour};
  \item inventory the cryptographic surfaces of a layer 2 by \emph{who can fix each one};
  \item explain why a hash-based or symmetric layer does not make a system post-quantum, and what to ask instead;
  \item describe Ethereum's exposure and its two-front plan: an emergency hard fork and account-abstraction-based migration;
  \item explain why an EdDSA chain admits a migration route that an ECDSA chain does not, and why generating ECDSA keys from a seed does not close the gap;
  \item bound a migration by the \emph{throughput} of the chain itself, and compute that bound;
  \item weigh the central controversy -- whether to freeze quantum-vulnerable coins -- and the honest uncertainty in the timeline.
\end{itemize}

\section{The blockchain threat model}
The immediate migration pressure is on signatures and pairing-based public-key systems.
Shor's algorithm (Chapter~\ref{ch:quantum}) threatens Bitcoin's ECDSA and Schnorr,
Ethereum's account ECDSA and consensus BLS signatures, and the public-key assumptions
behind SNARK and KZG systems~\cite{quantumhorizon2026,ethpqroadmap}. Grover gives
hashes a square-root rather than exponential weakening, but hashes still require
property- and output-length-specific assessment: HASH160, for example, has only a
nominal $2^{80}$ quantum preimage margin.

The subtlety that governs everything is \textbf{when a public key becomes visible}.
P2PK and P2TR expose a public key or tweaked public key when a UTXO is created;
P2PKH and P2WPKH usually reveal it on first spend, and reuse leaves later outputs
exposed. Other scripts must be assessed by their actual path and witness.

\begin{figure}[H]
\centering
\begin{tikzpicture}[node distance=6mm and 12mm, every node/.style={font=\scriptsize}]
  \node[flownode, align=center] (funded) {funded, never-spent address\\[-2pt]\scriptsize only a \emph{hash} of the public key is on chain};
  \node[flownode, align=center, right=of funded] (spend) {first spend\\[-2pt]\scriptsize the signature reveals the public key};
  \node[keynode, align=center, right=of spend] (exposed) {public key now on chain\\[-2pt]\scriptsize Shor can derive the private key};
  \draw[flowarrow] (funded) -- (spend);
  \draw[flowarrow] (spend) -- (exposed);
  \node[pqcteal] at ([yshift=-7mm]funded.south) {Grover only: safe};
  \node[pqcorange] at ([yshift=-7mm]exposed.south) {Shor: exposed};
\end{tikzpicture}
\caption{The public-key-exposure lifecycle. While an address has only ever received funds, its public key is hidden behind a hash and only Grover applies. The first spend publishes the public key permanently, exposing the account to Shor's algorithm. This is why exposed coins are considered at risk indefinitely: the ledger never forgets, so an attacker with a future quantum computer needs no fresh interception -- the target is already recorded.}
\label{fig:pubkey-exposure}
\end{figure}

\section{Bitcoin}
\subsection{What is exposed}
Bitcoin's exposure is precisely a function of output type. Coins in \emph{pay-to-public-key} (P2PK, the earliest, Satoshi-era outputs), raw multisig, and Taproot key-path outputs carry a public key on chain permanently and are long-exposure vulnerable. Coins in the hash-wrapped types (P2PKH, P2SH, P2WPKH) stay protected \emph{until the address is reused or spent}, at which point the key is revealed~\cite{bip360,chaincodepq2025}. BIP-361 states that, as of 1~March 2026, over \textbf{34\% of all bitcoin have revealed a public key on chain}~\cite{bip361}; an independent Chaincode Labs analysis puts the exposed fraction in the range of roughly one-fifth to one-half of supply~\cite{chaincodepq2025}. (Older and widely repeated figures attributing a specific amount to Satoshi-era P2PK are best treated as estimates rather than established fact.)

\subsection{Two kinds of exposure}
\label{sec:two-exposures}

The aggregate percentage hides the more useful structure. Published on-chain
analyses that break the exposed supply down by \emph{cause} agree on the
qualitative shape, even where their totals differ: roughly one quarter to one
third of it comes from output types that publish a key by design (P2PK, raw
multisig, Taproot key-path), and the clear majority --- on the order of
seven-tenths --- comes from ordinary hash-wrapped addresses that were
\textbf{reused after being spent from}. Figure~\ref{fig:pubkey-exposure} draws
the mechanism; the split tells us how often each branch is taken.

The distinction matters because the two causes have different remedies and
different politics:

\begin{itemize}
  \item \textbf{Design exposure} is a property of the protocol. No user action
        can undo it --- a P2PK output has its public key in the ledger forever,
        and the owner cannot retract it. Only a migration to a new output type
        helps, and for P2PK outputs whose keys are lost, nothing helps at all.
  \item \textbf{Behavioural exposure} is a property of how wallets were used. It
        was avoidable in principle, it is still accumulating today with every
        reused address, and --- crucially --- it is the part that education,
        wallet defaults, and fee policy can actually reduce \emph{before} any
        consensus change ships.
\end{itemize}

The practical consequence is that the largest single share of Bitcoin's quantum
exposure does not require a soft fork to start shrinking. It requires wallets to
stop reusing addresses, which is a change nobody has to be persuaded to ratify.
Any exposure figure quoted without this breakdown makes the problem look more
uniformly intractable than it is.

\begin{remark}[Reading exposure statistics]
On-chain exposure figures are analyst-dependent. They vary with the treatment of
dust and provably-unspendable outputs, with whether an address is counted by
balance or by output count, and with how Taproot outputs are classified. Quote a
range and a source, as this section does, rather than a single decimal figure.
\end{remark}

\subsection{The draft proposals}
Two Bitcoin Improvement Proposals anchor the response; both are \textbf{Draft and not activated} as of mid-2026:

\begin{itemize}
  \item \textbf{BIP-360, ``Pay-to-Merkle-Root'' (P2MR)}~\cite{bip360} -- formerly circulated as ``Pay-to-Quantum-Resistant-Hash'' (P2QRH) -- introduces a new output type, proposed as a backward-compatible soft fork. In its current form it is essentially Taproot with the key-path spend removed, so that no bare public key is ever committed; notably, the current draft \emph{defers} the choice of the actual post-quantum signature scheme (naming ML-DSA and SLH-DSA only as examples) to a future proposal, rather than mandating one.
  \item \textbf{BIP-361, ``Post-Quantum Migration and Legacy Signature Sunset''}~\cite{bip361} -- the migration-and-deadline layer. It proposes a phased sunset of ECDSA/Schnorr: a first phase (on the order of three years after activation) that stops sending funds to legacy address types, and a later flag-day (around five years) that restricts legacy signatures, so that only a holder who can prove ownership by a quantum-safe route can still move the coins.
\end{itemize}

\subsection{Choice of signature scheme}
\label{sec:bitcoin-scheme-choice}

BIP-360 deliberately leaves the scheme unspecified, so the question is open
rather than settled; this subsection summarizes what the evidence constrains. A 2026 Blockstream Research review of the lattice
candidates~\cite{blockstreamlattice2026} works the problem through for Bitcoin
specifically. Its central framing: for a chain, the cost metric is not the
signature size but the \textbf{combined public-key and signature size}, because
a spend writes both to the ledger. Under this metric the ranking of the schemes
changes.

\begin{table}[H]
\centering
\begin{tabular}{lrrl}
\toprule
Scheme & NIST level & $|pk| + |\sigma|$ & Note \\
\midrule
Schnorr (\texttt{secp256k1}) & --- & $96$ B & the incumbent; Shor-breakable \\
\midrule
FN-DSA-512 (Falcon)   & 1 & $1563$ B & smallest standardized \\
FN-DSA-1024 (Falcon)  & 5 & $3073$ B & \\
ML-DSA-44             & 2 & $3732$ B & \\
ML-DSA-65             & 3 & $5261$ B & \\
\midrule
SLH-DSA-128s          & 1 & $7888$ B & most conservative assumption \\
PORS+FP               & 1 & $3472$ B & non-standard; $2^{30}$ signing budget \\
\bottomrule
\end{tabular}
\caption{On-chain footprint of a single spend, following the comparison
in~\cite{blockstreamlattice2026}. Every post-quantum row is one to two orders of
magnitude above the incumbent, which is the quantity
Section~\ref{sec:migration-throughput} turns into a migration schedule.}
\label{tab:bitcoin-scheme-sizes}
\end{table}

Three consequences follow from Table~\ref{tab:bitcoin-scheme-sizes}.

\paragraph{Security level for long-lived outputs.}
A Bitcoin output may remain unspent for decades, so the relevant security
horizon is the lifetime of the output rather than of the software that created
it. If cryptanalytic progress later pushes a scheme below its intended
category, coins locked under the weakened keys remain exposed, and the freeze
question of Section~\ref{sec:freeze} returns. Large deployments already avoid
Category 1 lattice parameters on these grounds: Apple's iMessage PQ3 uses
Categories 3 and 5, and Cloudflare pairs X25519 with ML-KEM-768 rather than
ML-KEM-512, in both cases citing a margin against cryptanalytic progress over
the lifetime of the system. Hash-based schemes are commonly deployed at
Category 1, because the underlying assumption is older and better studied than
the lattice assumptions.

\paragraph{The cost of a security margin.}
Falcon-1024, at Category 5, has a smaller combined footprint (3073 B) than the
best stateless hash-based construction at Category 1 (7888 B) and than
ML-DSA-65, the standardized form of Dilithium~\cite{dilithium2018}
(Chapter~\ref{ch:mldsa}), at Category 3 (5261 B). For ML-DSA, raising the
category from 2 to 3 adds about 1.5 kB per spend. Falcon is therefore the only
candidate whose highest security category remains smaller than every
alternative at lower categories.

\paragraph{Absence of an intermediate parameter set.}
The recursive sampler of Chapter~\ref{ch:fndsa} halves the ring degree at every
step, so the admissible degrees are powers of two, and the parameter sets pass
from $n=512$ (approximately Category 1) directly to $n=1024$ (above
Category 5). There is no intermediate ring and therefore no intermediate
parameter set; a deployment must choose one of the two.

\paragraph{Assessment.}
The review~\cite{blockstreamlattice2026} recommends a staged transition. For
immediate deployment it recommends a hash-based scheme, on two grounds: the
hash assumption is the most mature of the alternatives, and FN-DSA is not a
finished standard (FIPS 206 was still in development at this book's snapshot,
so there is no fixed specification, no test vectors, and no audited library
ecosystem to integrate against). Once Falcon completes standardization, its
size advantage over every hash-based alternative changes the comparison. The
two families can also be combined: a stateful hash-based scheme with short
common-case signatures can use Falcon as its stateless fallback, at a much
smaller size than an SLH-DSA fallback for the rare recovery path.

\begin{remark}[Sizes alone do not decide the choice]
The smallest entry in Table~\ref{tab:bitcoin-scheme-sizes} does not by itself
determine the recommendation. Falcon has the smallest footprint but also the
least settled standardization status, the most delicate implementation
(Section~\ref{sec:falcon-hard}), and no practical wallet derivation
(Section~\ref{sec:pq-wallet-derivation}); these are the grounds on which the
review recommends the hash-based family for immediate deployment.
\end{remark}

\subsection{The wallet-derivation gap}
\label{sec:pq-wallet-derivation}

There is a second constraint on the choice, and it is not a size problem at all.
Bitcoin wallets are hierarchical and deterministic: from one master extended
public key, BIP-32~\cite{bip32} derives an unbounded supply of child public keys
using \emph{public data only}, with the secret key never involved. That is
what lets a watch-only hot wallet hand out fresh addresses while the seed stays
in cold storage, and it is the backbone of essentially every wallet in use.

\textbf{None of the standardized post-quantum signature schemes supports
it.}~\cite{blockstreamlattice2026} BIP-32's unhardened derivation is an addition
in the elliptic-curve group --- child key $=$ parent key $+$ tweak, on both the
secret and public sides --- and lattice schemes do not offer the same structure
for free. The known replacements are research constructions, and expensive ones:

\begin{itemize}
  \item \textbf{Dilithium} is the furthest along. Rerandomizable-key variants
        derive a child from a public update value and keep the two keys in
        step, at a combined footprint of roughly 6.5 to 8.3 kB depending on
        whether the wallet publishes an uncompressed key or carries a
        correction in every signature. The obstruction is that ML-DSA's
        public-key compression is not additive, so an update cannot simply be
        added to a compressed parent.
  \item \textbf{Falcon} has one published construction, which rerandomizes the
        lattice basis. Doing so inflates the signature-norm bound and pushes an
        on-chain signature to roughly \textbf{23.7 kB}, about $36\times$ a
        Falcon-512 signature --- and its stated parameters do not satisfy the
        scheme's own security condition, so a corrected version is larger
        still. There is currently no workable public-key derivation for Falcon.
\end{itemize}

The Dilithium variants also carry a deployment assumption that inverts one of
ML-DSA's own design decisions. Unlinkability of
derived addresses requires the matrix $A$ to be shared \emph{network-wide}. That
is formally sound under Module-LWE, but Chapter~\ref{ch:mldsa} derives $A$ from
a per-key seed precisely to rule out a backdoored matrix and to stop one
precomputation from applying to many targets. The preprocessing in dual-lattice
attacks depends only on the matrix, so a single network-wide $A$ would let one
large reduction effort apply to every key on the chain at once. At Dilithium-3
dimensions that attack remains far out of reach --- but the convention
concentrates the whole network's security on one instance, which is a different
risk profile from the one FIPS 204 chose.

\begin{remark}[What this costs a migration]
Section~\ref{sec:migration-throughput} treats the per-output cost $w$ as the
term a protocol can negotiate. Wallet derivation is a reminder that $w$ is not
the only thing at stake: a scheme with no unhardened derivation forces either a
different wallet architecture (every address pre-derived from the cold key and
stored) or a scheme variant that is not the standardized one. Neither is fatal,
and neither is free, and both are easy to omit from a comparison that stops at
signature size.
\end{remark}

A recurring safety mechanism is \emph{commit--reveal}: a holder first publishes a
binding commitment and only later reveals it. It is designed to mitigate front-running
under the stated confirmation-depth, reorganization, fee, timeout, and recovery
assumptions; it is not an unconditional guarantee.

\subsection{The freeze controversy}
\label{sec:freeze}
BIP-361's sunset implies that coins whose owners never migrate -- and coins for which no ownership proof is possible, notably P2PK -- become effectively unspendable, i.e.\ \emph{frozen}. This is the sharpest open dispute in Bitcoin's roadmap. Proponents argue a pre-announced sunset is defensive: a quantum attacker would otherwise \emph{steal} those coins outright, whereas a freeze at least preserves them. Critics call any forced migration confiscatory, contrary to Bitcoin's property-rights ethos, and warn that freezing long-dormant coins (including Satoshi's) sets a dangerous precedent. This is a genuine, unresolved political disagreement, not a settled plan, and should be read as such.

\section{Ethereum}
\subsection{What is exposed}
Ethereum addresses reveal an effectively public key once an ECDSA transaction has been
sent. BLS validator signatures, pairing curves behind SNARK verifiers, and KZG
commitments all rely on public-key assumptions threatened by Shor~\cite{ethpqroadmap,quantumhorizon2026}.
A comparison of quantum cost for secp256k1 and BLS12-381 needs group-specific resource
estimates; the 381-bit base-field name alone is not such an estimate. KZG migration
must distinguish blob retention from commitment binding, proof verification, and
consensus dependencies; blob pruning alone does not make the risk mild.

\subsection{The two-front plan}
\label{sec:eth-two-front}
Ethereum's response has an emergency track and a structural track; nothing post-quantum has yet shipped to mainnet.

\paragraph{The quantum-emergency hard fork.}
Vitalik Buterin's proposal ``How to hard-fork to save most users' funds in a quantum emergency''~\cite{buterinquantum2024} is the contingency plan. On evidence of large-scale theft, the chain would revert to the last pre-theft block, freeze legacy ECDSA accounts, and let users recover funds through a new transaction type carrying a \emph{STARK proof} that the sender knows the hash preimage secret behind their address -- authenticating by knowledge of a hash preimage rather than by an ECDSA signature. Because most users' recovery secret was never exposed on chain, this protects the majority, at the cost of a wallet-software update. Buterin notes the infrastructure could in principle be built now.

\paragraph{Account abstraction as the migration vehicle.}
The structural plan is to let accounts choose their own signature scheme. \textbf{EIP-7702} (``Set Code for EOAs,'' \emph{Final}, shipped in the 2025 Pectra upgrade)~\cite{eip7702} lets an externally-owned account delegate to contract code, turning it into a smart account with arbitrary verification logic -- although 7702 authorizations are themselves secp256k1-signed, so it is a stepping stone, not itself a post-quantum migration. The dedicated vehicle is \textbf{EIP-8141} (``Frame Transaction,'' \emph{Draft}, 2026, co-authored by Buterin)~\cite{eip8141}, which explicitly describes native account abstraction as ``a native off-ramp \dots{} to post-quantum secure systems,'' supporting a slot for an arbitrary (custom, post-quantum) signature scheme alongside the classical ones. At the consensus layer, the plan is to replace BLS aggregation with a hash-based signature of the XMSS/Winternitz family (Chapter~\ref{ch:hashots}), aggregated inside a purpose-built zero-knowledge virtual machine. Both pieces exist as public prototypes: \texttt{leanVM}, which describes itself as a ``minimal hash-based zkVM, for a Post-Quantum Ethereum,'' and \texttt{leanSig}, a Rust implementation of the proposed signature scheme whose README states plainly that it is unaudited and not for production~\cite{leanvm}, and to move KZG commitments to STARK- or lattice-based alternatives; the Ethereum Foundation's stated assessment targets first-layer post-quantum infrastructure around 2029~\cite{ethpqroadmap}. In August 2026 the Foundation announced a revision of this stack's cryptographic base: the proof system is to be built over binary fields, so that the hash function throughout --- in the signatures, in the Merkle commitments, and in the zkVM's own commitment scheme --- can be a standard one (SHA-2 or the BLAKE family) rather than an arithmetization-oriented hash such as Poseidon, on which the earlier prototypes relied. The announced schedule is a production-grade \texttt{leanVM} in 2027 and consensus-layer deployment in 2028~\cite{efbinary2026}.

\subsection{On-chain verification, revisited}
The structural migration path may introduce post-quantum signature verification on
chain; the emergency path instead proposes a STARK proof of seed/preimage knowledge.
These are different authentication mechanisms. Their costs are scheme-, parameter-,
implementation-, chain-, and gas-schedule-dependent, so a benchmark must state those
details, including code commit and measurement date, before comparing gas figures.

\begin{table}[H]
\centering
\small
\begin{tabular}{@{}p{2.5cm}p{3.6cm}p{4.0cm}p{4.0cm}@{}}
\toprule
 & \textbf{Bitcoin} & \textbf{Ethereum} & \textbf{Cosmos} \\
\midrule
Hardest constraint & governance & signature size at validator scale & heterogeneity across chains \\
Fund-securing signature & ECDSA / Schnorr (secp256k1) & ECDSA (accounts); BLS12-381 (consensus) & ed25519 (consensus) \\
Exposed today & P2PK, reused P2PKH, Taproot key-path & any account that has sent a transaction & validator keys; account keys \\
PQ scheme chosen & \textbf{none yet} & hash-based (XMSS/Winternitz family) & \textbf{ML-DSA-65} (FIPS 204) \\
Migration vehicle & new output type (BIP-360) & account abstraction; zkVM aggregation & opt-in key type $+$ key rotation \\
Emergency plan & legacy-signature sunset (BIP-361) & emergency hard fork with STARK recovery & per-chain governance \\
Fork style & soft fork & protocol upgrades / emergency hard fork & per-chain, genesis parameter \\
Status (mid-2026) & Draft BIPs, not activated & prototypes, unaudited; PQ infra $\sim$2029 & \textbf{shipped in SDK v0.55}, off by default \\
\bottomrule
\end{tabular}
\caption{Post-quantum migration at a glance, as of mid-2026. The three ecosystems share one physical fact -- post-quantum signatures are one to two orders of magnitude larger than what they replace -- but their binding constraints differ, and so do their mechanisms. Note that progress inverts exposure: Bitcoin, with the sharpest exposure, has specified no post-quantum signature scheme, while Cosmos has one in shipped code.}
\label{tab:chain-migration}
\end{table}


\section{Administrative keys: the concentrated case}
\label{sec:admin-keys}

Everything so far has treated the threat as distributed: many users, each holding
their own coins behind their own key, migrating at their own pace. One class of
target inverts that structure completely, and it is under-discussed relative to
its size.

A stablecoin is not a native asset but an upgradeable smart contract, and its
behaviour is governed by a handful of privileged roles held by ordinary
externally-owned accounts --- which is to say, by ordinary secp256k1 keys with
ordinary on-chain public keys. The roles vary by issuer, but the pattern is
consistent:

\begin{itemize}
  \item a \textbf{proxy admin}, which can replace the contract's implementation
        and therefore its entire logic;
  \item a \textbf{minter} or master-minter role, which can issue tokens without
        corresponding reserves;
  \item a \textbf{blacklister} or freezer, which can immobilize balances;
  \item a \textbf{pauser}, which can halt all transfers.
\end{itemize}

The security properties here are unlike the rest of this chapter in three ways.
The number of keys is tiny, so an attacker needs a handful of Shor executions
rather than a campaign. The value controlled is enormous and shared --- compromise
does not drain one wallet but affects every holder of the token at once, along
with every lending protocol, bridge, and exchange that treats it as collateral.
And the damage is not bounded by the balance, because a proxy admin can
\emph{rewrite the accounting rules} rather than merely move funds.

Two mitigations that work elsewhere do not help. Multisignature control divides
the private keys among several holders, but a quantum adversary derives each one
independently from its published public key, so an $m$-of-$n$ scheme raises the
attacker's cost by a factor of $m$, not exponentially --- the same linear
degradation noted for user wallets. Hardware custody protects a key against
extraction from a device, and Shor does not extract keys from devices; it derives
them from the ledger.

The redeeming feature is that this is the one part of the problem with a
\emph{responsible party}. Nobody can compel millions of users to rotate their
wallets, but a stablecoin issuer can migrate its own admin roles to post-quantum
authorization on its own schedule --- provided the chain gives it a post-quantum
verification path to migrate \emph{to}, which returns the problem to the protocol
roadmaps above. Concentrated control is the reason this exposure is dangerous,
and also the reason it is the most tractable item in the chapter.

\section{Cosmos}

Bitcoin and Ethereum are the chains the quantum discussion usually names, and on
both of them the post-quantum work is still a proposal. The ecosystem that has
gone furthest in shipped code is neither of them.

\subsection{An opt-in post-quantum consensus key}
Cosmos SDK v0.55 registers \textbf{ML-DSA-65} (FIPS~204, Chapter~\ref{ch:mldsa})
as a supported validator consensus key type~\cite{cosmossdkpq,cosmospqdocs}. This is code in
the tree, not a draft: the \texttt{cosmos.crypto.mldsa65} proto package, the
Amino routes, and an \texttt{x/auth} signature-verification cost parameter all
ship enabled, and the \texttt{init} and \texttt{testnet} commands accept
\texttt{-{}-consensus-key-algo ml\_dsa\_65}.

What ships enabled is the \emph{capability}, not the change.
\texttt{genesis.consensus\_params.validator.pub\_key\_types} still defaults to
\texttt{["ed25519"]}, so an existing chain is unaffected until its own
governance adds the new type. Because validator consensus key \emph{rotation}
ships in the same release, an existing chain can adopt post-quantum keys by
adding the type and having validators rotate, rather than by starting a new
genesis. That is a materially gentler migration path than a flag day.

\subsection{What the size actually costs}
The operational consequences are stated plainly, and they are the same physical
fact that shapes every roadmap in this chapter: an ML-DSA-65 public key is
\textbf{1952 bytes} against ed25519's 32, and a signature is \textbf{3309 bytes}
against 64 -- roughly $60\times$ and $50\times$. CometBFT's maximum signature
size and per-validator commit-signature budget had to be raised to fit them, and
chains enabling the type are told to revisit their block size and gossip framing
limits. Every validator signs every block, so this lands squarely on the
consensus hot path.

These two numbers are exactly what FIPS~204 prescribes for ML-DSA-65, and the
reader can confirm them rather than take them on trust: the implementation of
Chapter~\ref{ch:mldsa} derives both from the parameter set,
\[
|pk| = 32 + 32k(\mathrm{bitlen}(q-1) - d) = 1952,
\]
with the signature length following from $\lambda$, $\ell$, $\gamma_1$, $\omega$
and $k$ in the same way.

\subsection{The interoperability hazard}
Cosmos surfaces a migration problem that a single monolithic chain never
encounters, and it is the most transferable lesson in this chapter. IBC light
clients on a \emph{counterparty} chain verify your validator set's commit
signatures using the counterparty's own compiled-in cryptography. If validators
holding sufficient voting power start signing with a key type a counterparty
cannot verify, headers fail verification there, packet flow with that chain
stops, and the light client eventually expires~\cite{cosmossdkpq}. The
counterparty needs only the verification code, not the key type itself -- but it
does need it, \emph{first}.

In an interoperable ecosystem, then, post-quantum migration is a coordination
problem before it is a cryptographic one, and its failure mode is quiet: nothing
breaks at the moment of upgrade, only later, when a connection that used to
carry packets stops doing so. Any multi-chain migration plan needs a sequencing
story, and no published one yet covers a large IBC graph.


\section{Seeds, preimages, and the EdDSA asymmetry}
\label{sec:eddsa-asymmetry}

Every migration route so far has required the owner to act \emph{before} the
attacker does. There is one family of mechanisms that does not, and it turns on a
detail of key generation that most readers will never have had reason to notice.

\subsection{The recovery problem}
Suppose the worst has already happened: a quantum adversary has run Shor's
algorithm against an exposed public key and holds a valid private key. The chain
can protect the coins only by refusing legacy signatures --- but then the
\emph{honest owner} cannot spend either, because their private key is exactly the
one the attacker now has. Recovering the coins requires the owner to demonstrate
something the attacker cannot: not possession of the private key, which is no
longer secret, but possession of some secret \emph{upstream} of it.

Ethereum's emergency plan (Section~\ref{sec:eth-two-front}) is precisely this
move: authenticate by proving knowledge of a hash preimage rather than by signing.
What makes it work is that the recovery secret was never published, so Shor never
touches it. What makes it \emph{ad hoc} is that Ethereum must arrange for such a
secret to exist.

\subsection{RFC 8032 arranges it by construction}
EdDSA does not need the arrangement, because its key generation already has this
shape. Under RFC~8032~\cite{rfc8032} an Ed25519 key begins as a 32-byte
\emph{seed} $k$, and the signing scalar is derived by hashing it:
\[
  h = \mathrm{SHA\text{-}512}(k), \qquad
  s = \mathrm{clamp}\big(h[0..31]\big), \qquad
  A = s\cdot B ,
\]
with the upper half $h[32..63]$ retained as the deterministic nonce prefix. The
value stored by the wallet is $k$; the value that the public key commits to is
$s$.

This one hash is the whole asymmetry. Shor's algorithm computes discrete
logarithms, so from $A$ an attacker recovers $s$ --- and stops there, because
recovering $k$ from $s$ is a preimage search on SHA-512, which Grover reduces
only quadratically. The honest owner knows $k$ and the attacker does not, even
after the curve is completely broken. The owner can therefore prove, in zero
knowledge,
\[
  \exists\, k \;:\; \mathrm{clamp}\big(\mathrm{SHA\text{-}512}(k)[0..31]\big)\cdot B = A ,
\]
binding the existing address $A$ to a fresh post-quantum key, with no need to
have moved the funds beforehand~\cite{eddsapq2025}.

The distinction turns on the last clause:

\begin{itemize}
  \item A \textbf{pre-commitment} scheme --- publish a commitment to a future
        post-quantum key, as BIP-360's commit--reveal does --- requires every
        owner to \emph{act before Q-Day}.
  \item A \textbf{seed-preimage proof} lets an owner recover \emph{after Q-Day,
        having done nothing at all}.
\end{itemize}

Given that SegWit needed over three years to reach 80\% of transactions while
offering an immediate fee discount, ``requires no prior user action'' is not a
minor convenience. It is the difference between a mechanism that reaches dormant
coins and one that does not.

\subsection{Why not simply derive ECDSA keys from a seed?}
\label{sec:why-not-ecdsa-seed}

The obvious question is why ECDSA chains cannot adopt the same trick: define
$sk = H(\text{seed}) \bmod n$, keep the seed, and prove the preimage later.
Cryptographically, nothing prevents it. ECDSA places no constraint on how the
scalar is produced, so the construction is sound and the proof statement is the
same shape.

More than that, \textbf{it has already been done}. BIP-32 hierarchical
deterministic wallets~\cite{bip32} derive every signing key from a master seed
through a chain of HMAC-SHA512 steps, and BIP-39 puts a mnemonic and PBKDF2 in
front of that. Most Bitcoin and Ethereum keys in existence today do have a hash
between the secret their owner wrote down and the scalar their address commits
to. The blanket claim that ECDSA chains lack this layer is a statement about the
\emph{protocol}, and it is too strong as a statement about \emph{wallets}.

The gap is nevertheless real, and it is instructive because it is not a
cryptographic gap at all. Three things separate a protocol invariant from a
widespread convention:

\begin{enumerate}
  \item \textbf{Universality.} RFC~8032 makes seed-then-hash the \emph{definition}
        of Ed25519 key generation, so every conforming key has the structure. In
        ECDSA it is a wallet choice, and the exceptions are exactly the
        high-value historical cases: early Bitcoin keys from \texttt{dumpprivkey},
        brainwallets, paper wallets, exchange- and HSM-generated keys, and
        anything produced before BIP-32 existed in 2012.
  \item \textbf{Detectability.} Given an address, nothing on chain reveals
        whether it has BIP-32 ancestry. A consensus rule cannot offer a recovery
        path to seed-derived keys and freeze the rest, because it cannot tell them
        apart in advance; owners of raw-entropy keys would discover the problem
        only at the moment they tried to recover.
  \item \textbf{A single statement to prove.} Ed25519 needs one circuit. ECDSA
        recovery would have to cover BIP-32 against SLIP-0010 against Electrum's
        scheme against raw entropy, hardened against non-hardened derivation, and
        an arbitrary derivation path --- either as a much larger circuit or by
        revealing the path as a public input, which is probably acceptable but is
        one more thing to specify and get adopted.
\end{enumerate}

And for keys generated \emph{from now on}, the idea dissolves entirely. Changing
key generation requires a wallet software update; if a wallet is being updated
anyway, it should generate an ML-DSA key rather than a cleverly-derived ECDSA
one. The seed-preimage route has value only \emph{retroactively}, for coins
already sitting at addresses whose owners cannot or will not act --- which is
also, not coincidentally, the population every other mechanism in this chapter
fails to reach.

\begin{table}[H]
\centering
\small
\begin{tabular}{@{}p{3.4cm}p{5.0cm}p{5.4cm}@{}}
\toprule
 & \textbf{EdDSA chains} & \textbf{ECDSA chains} \\
\midrule
Seed-to-scalar hash & mandated by RFC~8032 & wallet convention (BIP-32/39) \\
Coverage & every conforming key & most, but not all, and not the oldest \\
Detectable on chain & yes, by construction & \textbf{no} \\
Derivation to prove & one function & many schemes, plus a path \\
Usable as a consensus rule & plausibly & not without classifying keys it cannot classify \\
\bottomrule
\end{tabular}
\caption{Why the same cryptographic idea supports a chain-wide policy in one case
and not the other. The difference is not the hash; it is whether the hash is a
protocol invariant or a wallet habit.}
\label{tab:eddsa-ecdsa}
\end{table}

\subsection{Caveats that apply to both}
Three conditions are easy to overlook, and each one connects to material
elsewhere in this chapter.

\paragraph{It is recovery, not protection.}
A seed proof does nothing while legacy signatures are still valid: the attacker
holds a working private key and simply spends with it, faster than the owner will
notice. The mechanism only has force on a chain that has \emph{already disabled}
legacy authorization. It therefore does not sidestep the freeze controversy of
Section~\ref{sec:freeze} --- it presupposes that the sunset argument has been won.

\paragraph{The proof system must itself be post-quantum.}
Proving a hash preimage inside a Groth16 verifier over a pairing-friendly curve
would produce an argument whose soundness Shor breaks, and the resulting system
would be no better than the ECDSA it replaced. This is exactly the error
catalogued in Section~\ref{sec:symmetric-plumbing}: a hash-based
\emph{statement} carried by an elliptic-curve \emph{argument}. A STARK or other
hash-based proof system is required, and it must be checked all the way down,
including its lookup and memory arguments.

\paragraph{The proof must be bound to the destination.}
Recovery proofs are published. If a proof establishes only ``someone knows the
seed behind $A$'' without committing to where the funds go, an attacker watching
the mempool can replay it in a transaction paying themselves. The proof statement
must include the destination, which makes it a signature over the recovery
transaction rather than a free-floating attestation --- the same front-running
hazard, with the same class of fix, as the commit--reveal construction discussed
above.

\begin{remark}[BIP-32 degrades separately under Shor]
Non-hardened BIP-32 derivation sets $k_{\text{child}} = k_{\text{parent}} + I_L$
where $I_L$ is computable from the parent \emph{public} key and chain code. An
account-level extended public key is routinely shared with watch-only wallets and
light-client servers; against Shor, that exposure yields the account scalar and
hence every non-hardened child. The master seed stays safe, so the recovery
argument above still holds --- but the cost of a single exposure is an entire
account branch, not one address.
\end{remark}

\section{Layer 2s: the composition problem}
\label{sec:l2-composition}

The three chains above each migrate one system. A \emph{layer 2} does not have
that luxury. A rollup or sidechain that settles to Bitcoin is a composition of
systems designed separately, migrating on separate timelines, owned by separate
parties: it settles to a base layer it cannot change, usually borrows a
consensus stack from a third ecosystem, and operates a bridge whose trust root
is cryptography of its own choosing~\cite{btcl2pqc}.

That composition changes the question. For a single chain one may sensibly ask
``when will it be post-quantum?'' For a layer 2 the only answerable question is
\emph{which surfaces do we own, and what is the residual risk on the rest?}
Table~\ref{tab:l2-surfaces} is the inventory that question needs.

\begin{table}[H]
\centering
\small
\begin{tabular}{@{}p{3.6cm}p{4.3cm}p{2.0cm}p{3.6cm}@{}}
\toprule
Surface & Typical scheme & Shor? & Who can fix it \\
\midrule
L1 settlement outputs & secp256k1 / Schnorr & yes & \textbf{base layer only} \\
Bridge operator keys & secp256k1, Schnorr, BLS & yes & the L2 \\
Bridge proof system & Groth16/PLONK, or STARK & \emph{check inside} & the L2, often with upstream \\
Bridge commitments & Merkle, Winternitz & \textbf{no} & already safe \\
L2 consensus keys & ed25519, BLS12-381 & yes & the L2's consensus stack \\
Contract admin / upgrade keys & secp256k1 (often multisig) & yes & the role holder, unilaterally \\
EVM crypto precompiles & \texttt{ecrecover}, BN254 pairing, KZG & yes & \textbf{upstream; unmigratable} \\
\bottomrule
\end{tabular}
\caption{Cryptographic surfaces of a Bitcoin layer 2, keyed on \emph{who can fix each one}. The ownership column, not the scheme column, is what determines a migration plan.}
\label{tab:l2-surfaces}
\end{table}

\subsection{The reach of a symmetric layer}
\label{sec:symmetric-plumbing}
The most useful lesson from auditing such a stack is a negative one, and it
generalizes well beyond blockchains: \textbf{a hash-based or symmetric layer
protects only what it carries.}

Three instances of the same mistake, all from one real bridge:
\begin{itemize}
  \item Bitcoin-side bit commitments implemented as \textbf{Winternitz one-time
        signatures} (Chapter~\ref{ch:hashots}) --- hash-based, and genuinely
        post-quantum --- were committing the elements of a \emph{Groth16 proof
        over a pairing-friendly curve}. The commitment survives Shor; the thing
        committed does not.
  \item A \textbf{garbled-circuit} evaluator, built from AES and a hash
        function and therefore symmetric throughout, was garbling a
        \emph{Groth16 verifier}. Garbling changes how the verifier executes, not
        what it asserts.
  \item A \textbf{FRI-based STARK}, whose polynomial commitment is hash-based,
        used an \emph{elliptic-curve multiset hash} for its offline
        memory-consistency argument. A quantum adversary never attacks FRI; it
        forges the memory check~\cite{ziren276}.
\end{itemize}

In each case a correct statement --- ``this layer is hash-based'' --- supported a
false conclusion. The habit to build is to ask what a layer \emph{carries}, and
to keep asking until the answer is a hardness assumption rather than a
primitive.

\subsection{The several assumptions inside one proof system}
The third instance deserves emphasis because it is the least obvious. A proof
system is not one assumption but several, and the headline commitment scheme is
only the most visible. Permutation arguments, lookup arguments and offline
memory checks are all places where an elliptic-curve accumulator can sit inside
an otherwise hash-based system.

The concrete example is a zkVM whose multiset hash mapped memory values to
curve points, so that soundness of the memory argument rested explicitly on
discrete log --- documented by the project itself, with the fix being a
lattice-based multiset hash (LtHash) in the sense of~\cite{zisklthash}, and a
working prototype~\cite{ziren276}. The encouraging part of that story is that
the fix is a \emph{primitive swap inside an existing argument}, not an
architectural change. The cautionary part is that the exposure was visible only
in the project's own issue tracker, not in its dependency list.

\subsection{The exposure that cannot be migrated}
One row of Table~\ref{tab:l2-surfaces} behaves unlike the others. An account can
adopt a new signature scheme because its owner opts in. A \emph{precompile}
cannot: on an EVM chain the elliptic-curve and pairing precompiles
(\texttt{ecrecover}, the BN254 group and pairing operations, the BLS12-381 suite,
KZG point evaluation) are consensus rules, and the contracts that call them are
immutable bytecode. Every on-chain SNARK verifier calls the pairing precompile;
that code cannot be upgraded, and generally nobody has the authority to upgrade
it.

So this exposure is fixed at deployment time. After a chain is live, the only
available action is an \emph{inventory} --- which deployed contracts depend on
which precompile, and which of those are consensus- or bridge-critical --- and
that inventory is worth doing early, because an immutable caller can be
identified now but never migrated later.

\subsection{Replacing signature aggregation}
A bridge or consensus layer that aggregates signatures gets a large constant
factor from BLS: many signatures verify as one. No standardized post-quantum
signature aggregates. ML-DSA and SLH-DSA have no aggregation, so a naive
replacement turns one signature into $N$, each one to two orders of magnitude
larger than what it replaced --- for a committee of twenty ML-DSA-65 signers,
tens of kilobytes where there were tens of bytes.

This is the same wall Ethereum hit at the consensus layer, and the reason its
answer is a hash-based signature scheme aggregated inside a zkVM rather than a
drop-in replacement (Section~\ref{sec:eth-two-front}). Any system whose design
depends on aggregation should expect its post-quantum migration to be an
architectural project, not a key-type change.

The word \emph{aggregation} covers three different requirements, and the
post-quantum options differ for each.

\paragraph{Recovering aggregation with a proof system.}
The Ethereum route replaces the aggregate with a proof: each validator signs
with a hash-based scheme, and a zkVM proves that $N$ signatures verified. It
recovers a constant-size object for any $N$ and any set of messages, at the cost
of an entire proving pipeline, and it is the only route with a credible path to
$N$ in the hundreds of thousands.

\paragraph{Threshold signing.}
Most bridge and committee designs need less than general aggregation: many
signers, \emph{one} message, accepted once a threshold of them has signed. That
is the shape a \textbf{threshold signature} has natively --- $t$ of $n$ parties
run a protocol and jointly emit a single \emph{ordinary} signature under a joint
public key --- and unlike aggregation it leaves the verifier untouched: the
chain keeps checking one FIPS~204 signature with a stock verifier, whatever the
size of the signing set. Two 2026 constructions make this concrete for ML-DSA
and differ in the dimension that matters for a committee.
\textbf{Quorus}~\cite{quorus2025} supports any $t$-of-$n$ by evaluating the
ML-DSA signing algorithm inside a multi-party computation, at roughly
$100$\,kB of communication per party per rejection-sampling round;
\textbf{Celi et al.}~\cite{thresholdmldsa2026} give a lighter protocol limited to
about six parties at up to 1\,MB per party. Both produce signatures and keys
byte-identical to single-signer ML-DSA. For a signing set in the tens only the
first scales.

The price is a change in \emph{signing shape} rather than in key type. Today
each committee member signs independently and offline and a bitmap records who
did; threshold signing replaces that with a live multi-round protocol among the
selected quorum, so a liveness-tolerant design acquires a signing-time
availability requirement. And none of it is standardized: NIST's first call for
multi-party threshold schemes, IR~8214C~\cite{nistthreshold}, was finalized only
in January 2026. For a peg or bridge trust root, depending on a construction
with no FIPS number is a governance decision as much as a technical one, which
argues for migrating any key that needs no aggregation first.

\paragraph{Post-hoc aggregation.}
The option that sounds best --- compress $N$ existing signatures into one small
object without changing how anyone signs --- works worst. The first aggregate
signature for Fiat--Shamir-with-aborts schemes~\cite{boudgousttakahashi2023}
reports ``quite small compression rates'' in its authors' words, and aggregates
\emph{distinct} messages where a committee has one. The strong result in this
line, aggregating with LaBRADOR (Section~\ref{sec:pq-proof-systems}), applies
to Falcon and not to ML-DSA. If aggregability is a first-order requirement, it
is a consideration in \emph{which} post-quantum scheme to adopt, not a problem
to solve afterwards.

\subsection{The commitments inside the proof system}
\label{sec:pq-proof-systems}
The question of the preceding subsections --- what does each layer carry ---
applies to the proof system itself. There it becomes a question about the
\textbf{polynomial commitment}, because that is where a quantum adversary
attacks: the rest of a modern argument is field arithmetic over a transcript.

Two families of commitment are in use, and they differ in what they assume
rather than in how fast they run.

\begin{itemize}
  \item \textbf{Hash-based} commitments --- FRI, and its successors STIR and
        WHIR --- open a Merkle tree, so the assumption is a random oracle plus
        collision resistance, with nothing algebraic underneath. This is what
        makes a STARK post-quantum, and it is why the memory-check bug of
        Section~\ref{sec:symmetric-plumbing} was a bug rather than a design
        choice: an elliptic-curve accumulator had been placed inside a system
        whose commitment assumed nothing of the kind.
  \item \textbf{Lattice-based} commitments --- the LaBRADOR
        line~\cite{labrador2023} --- commit with an \textbf{Ajtai commitment},
        which is a matrix--vector product over the module of
        Chapter~\ref{ch:mlwe} and is opened by ring arithmetic rather than by
        hashing. There is no hash anywhere in the opening.
\end{itemize}

\subsubsection{The Ajtai commitment}
The whole scheme is one line of linear algebra. Fix the ring
$R_q=\mathbb{Z}_q[X]/(X^d+1)$ of Chapter~\ref{ch:mlwe} --- Greyhound's
implementation takes $d=64$ --- and a norm bound $\beta$.

The figures in this section follow the convention of the lattice zero-knowledge
literature~\cite{lnp2022} rather than the numbered pseudocode used elsewhere in
this book: a header declaring what is public and what is private, then the two
parties in columns with the messages between them, and the verifier's checks
gathered into a closing \emph{Accept iff}. For these protocols the interesting
content is who knows what and who sends when, and that convention shows it.

\begin{figure}[H]
\begin{protocol}
\textbf{Public information:} $\mathbf{A} \in R_q^{\kappa\times n}$ uniformly
random, norm bound $\beta$, commitment $\mathbf{t} \in R_q^{\kappa}$.\\
\textbf{Private information:} $\mathbf{s} \in R^{n}$ such that
\[
  \mathbf{t} = \mathbf{A}\mathbf{s} \bmod q
  \qquad\text{and}\qquad
  \|\mathbf{s}\| \le \beta .
\]
\protorule
\begin{tabular}{@{}p{0.30\linewidth}@{}p{0.30\linewidth}@{}p{0.36\linewidth}@{}}
\Prover & & \Verifier \\[0.9em]
 & \centering$\msgto{\mathbf{s}}$ & \\[0.6em]
 & & Accept iff \\
 & & \quad 1. $\|\mathbf{s}\| \le \beta$ \\
 & & \quad 2. $\mathbf{A}\mathbf{s} = \mathbf{t}$ \\
\end{tabular}
\end{protocol}
\caption{The Ajtai commitment and its opening. There is no trapdoor and no
structured reference string: $\mathbf{A}$ is public randomness, exactly as a
Merkle tree's hash function is. Note what is on the arrow --- the opening is the
entire witness, which is the deficiency the rest of this section is about.}
\label{fig:ajtai-commit}
\end{figure}

Three properties, each doing real work.

\paragraph{Compression.}
The commitment is $\kappa$ ring elements no matter how large $n$ is, so
$\kappa d\lceil\log_2 q\rceil$ bits stand for a witness of any length. A second
property accompanies this one: because the hardness of SIS does not depend on
the \emph{dimension} of the solution, lengthening the witness does not weaken
the assumption either, so a longer witness costs neither space nor security.
$\mathbf{A}$ is public randomness, with no trusted setup and no structured
reference string, exactly as for a Merkle root.

\paragraph{Binding.}
Suppose $\mathbf{s}\ne\mathbf{s}'$ both open $\mathbf{t}$. Subtracting the two
verification equations gives
\[
  \mathbf{A}(\mathbf{s}-\mathbf{s}') = \mathbf{0} \bmod q,
  \qquad
  \mathbf{s}-\mathbf{s}' \ne \mathbf{0},
  \qquad
  \|\mathbf{s}-\mathbf{s}'\| \le 2\beta ,
\]
which is precisely a solution to Module-SIS at bound $2\beta$.

Read as a reduction, this is tight and direct. An adversary against binding is
handed the M-SIS matrix unmodified (the challenge distribution and the
commitment key are the same uniform $\mathbf{A}$, so nothing has to be
simulated), and its output is converted into a solution by a single
subtraction. There
is no rewinding, no random oracle, and no loss:
$\mathrm{Adv}^{\mathrm{bind}} \le \mathrm{Adv}^{\text{M-SIS}}_{2\beta}$.

The comparison with a Merkle commitment cuts both ways. Binding here rests on
the hardness of one specific algebraic problem, Module-SIS in a fixed ring at a
fixed bound, where a Merkle commitment asks only that some hash function be
collision-resistant; in that sense the lattice assumption is the stronger one,
and structure is where the failures earlier in this section came from. Against
that, M-SIS is a falsifiable, parameterized assumption used through a tight
reduction, whereas a STARK's soundness argument goes through the random-oracle
model, an idealization of the hash rather than an assumption about it, together
with the proximity-gap conjectures of its proof of proximity. The lattice
commitment assumes more, and states precisely what it assumes; the hash-based
one assumes less, through a model that cannot be falsified.

\paragraph{Hiding.}
As drawn, Figure~\ref{fig:ajtai-commit} is a \emph{deterministic} function of
$\mathbf{s}$: it binds, but it emphatically does not hide, since testing a guess
means recomputing $\mathbf{A}\mathbf{s}$ and comparing. That stripped form is
what a succinct argument uses, for a reason returned to below. Stated in full,
the Ajtai scheme carries randomness $\mathbf{s}_2$ alongside the message
$\mathbf{s}_1$, both short:
\[
  \mathbf{A}_1\mathbf{s}_1 + \mathbf{A}_2\mathbf{s}_2 = \mathbf{t} \bmod q .
\]
Binding is the same subtraction as before. Hiding now follows because
$(\mathbf{A}_2,\mathbf{A}_2\mathbf{s}_2)$ is indistinguishable from uniform,
which is decisional Module-LWE (Chapter~\ref{ch:mlwe}). \textbf{Binding from
SIS, hiding from LWE} --- the two hard problems of Part~\ref{part:lattices}, one
each, on the same matrix.

The alternative is the \textbf{BDLOP} commitment, where only the randomness
$\mathbf{s}$ need be short and the message may be arbitrary modulo $q$:
\[
  \begin{bmatrix}\mathbf{A}\\ \mathbf{B}\end{bmatrix}\mathbf{s}
  \;+\;\begin{bmatrix}\mathbf{0}\\ \mathbf{m}\end{bmatrix}
  \;=\;\begin{bmatrix}\mathbf{t}_A\\ \mathbf{t}_B\end{bmatrix} \bmod q .
\]
The two schemes trade against each other in a way that determines the shape of
the protocols built on them. Ajtai demands a \emph{short} message, but
the message's dimension does not enlarge the commitment at all --- and since the
hardness of SIS does not really depend on the dimension of the solution, a longer
witness costs no security either. BDLOP accepts arbitrary messages and lets
further elements be appended for one ring element apiece, but its commitment
\emph{and} its opening grow linearly in the message. So one is for long witnesses
with small coefficients and the other for short witnesses with large ones, and
protocols need both: Lyubashevsky, Nguyen and Plan\c{c}on therefore fuse them
into a single ABDLOP commitment~\cite{lnp2022}, putting the long witness in the
Ajtai part and the low-dimensional auxiliary terms a protocol accumulates in the
BDLOP part, which yields the Ajtai commitment for free because one $\mathbf{t}_A$
serves both.

Both properties are computational, and that is forced rather than chosen: no
commitment is \emph{simultaneously} statistically binding and statistically
hiding, since an unbounded adversary breaks whichever was not paid for. Pushing
the randomness up to statistical hiding (a regularity argument over $R_q$) makes
it long, and the opening has to stay short for binding to bite, so the parameters
pull against each other.

Which is why the stripped form is the one that appears in this section. A proof
system needs the hiding half only if it wants zero knowledge, and these do not by
default --- LaBRADOR and Greyhound are arguments of \emph{knowledge}, with zero
knowledge available as a costed variant. Worth knowing before assuming that a
``commitment'' inside a lattice argument conceals anything: usually it is there
to bind a prover, not to keep a secret.

\paragraph{Linear homomorphism and norm growth.}
$\mathbf{A}\mathbf{s}_1+\mathbf{A}\mathbf{s}_2=\mathbf{A}(\mathbf{s}_1+\mathbf{s}_2)$,
which is why witnesses can be folded and batched at all. But the bound moves
with them: the sum opens to something of norm up to $2\beta$, a challenge-weighted
combination to more, and every round of folding must budget for the growth. In a
Merkle tree an opening is exact and composition is free. Here \textbf{norms are
the currency}, and the parameters are a norm budget.

\subsubsection{The opening problem}
Note what Figure~\ref{fig:ajtai-commit} does \emph{not} give: opening it means
sending $\mathbf{s}$, all $n$ ring elements of it. The commitment is short but
its opening is as long as the witness --- the exact opposite of a Merkle tree,
whose commitment is short \emph{and} whose openings are logarithmic, at the price
of a hash per level.

The central problem of lattice succinct arguments is therefore to convince a
verifier that a short opening exists without sending it. LaBRADOR's answer is
recursion --- reduce a system of quadratic relations
over $R_q$ to a smaller system of the same shape, commit, and repeat, so the
transcript is a sequence of commitments rather than a witness.

\subsubsection{The LaBRADOR argument}
\label{sec:labrador-outline}
LaBRADOR~\cite{labrador2023} is a proof of knowledge for one fixed shape of
statement, the \textbf{dot-product constraint system}: the witness is $r$ short
vectors $\mathbf{s}_1,\dots,\mathbf{s}_r\in R_q^{n}$ with
$\sum_i\|\mathbf{s}_i\|^2\le\beta^2$, and the constraints are polynomial
identities that are quadratic in the inner products,
\[
  \sum_{i,j} a_{ij}\,\langle \mathbf{s}_i,\mathbf{s}_j\rangle
  \;+\;\sum_i \langle \boldsymbol{\varphi}_i,\mathbf{s}_i\rangle
  \;+\; b \;=\; 0 \qquad\text{in } R_q ,
\]
one such equation per constraint. The public data of a constraint are of
three kinds. The $a_{ij}\in R_q$ are scalar coefficients on the quadratic
terms, the inner products of two witness vectors. Each
$\boldsymbol{\varphi}_i\in R_q^{n}$ is a public vector of the same length as
$\mathbf{s}_i$, so that $\langle\boldsymbol{\varphi}_i,\mathbf{s}_i\rangle$ is
the linear term in $\mathbf{s}_i$; the tuple
$(\boldsymbol{\varphi}_1,\dots,\boldsymbol{\varphi}_r)$ is the linear
coefficient vector of the constraint. And $b\in R_q$ is the constant. A
constraint is therefore a quadratic form in the witness with the inner product
as its only nonlinearity, and a purely linear constraint is the case in which
every $a_{ij}$ is zero.

Two examples show how a statement is written in this form. The opening of an
Ajtai commitment, $\mathbf{A}\mathbf{s}=\mathbf{t}$ with
$\mathbf{A}\in R_q^{\kappa\times n}$, is $\kappa$ linear constraints, the $k$-th
of them having $r=1$, $\boldsymbol{\varphi}_1$ equal to the $k$-th row of
$\mathbf{A}$, and $b=-t_k$; the shortness of $\mathbf{s}$ is not a constraint
but the norm bound $\beta$ carried alongside the system. A claim about an inner
product, such as $\langle\mathbf{s}_1,\mathbf{s}_2\rangle=v$, is one quadratic
constraint with $a_{12}=1$, no linear part, and $b=-v$. R1CS, and the
verification equations of lattice signatures, translate into collections of
such constraints in the same way; that is what makes LaBRADOR a
general-purpose argument rather than a proof of one relation.

The split into a quadratic part and a linear part is not cosmetic, because the
protocol treats the two differently. The quadratic terms are handled by
committing to the inner products $g_{ij}=\langle\mathbf{s}_i,\mathbf{s}_j\rangle$
themselves. The linear terms cannot be committed the same way, since after
amortization the verifier holds a single combination
$\mathbf{z}=\sum_ic_i\mathbf{s}_i$ and the cross terms
$\langle\boldsymbol{\varphi}_i,\mathbf{s}_j\rangle$ mix; the prover commits
the mixed values $h_{ij}$ and the verifier must form the combined coefficient
vector $\sum_ic_i\boldsymbol{\varphi}_i$ itself. Forming it touches every
$\boldsymbol{\varphi}_i$, and the $\boldsymbol{\varphi}_i$ together are as
large as the public statement, which is where the linear verifier of the
figure below comes from.

A single round of the protocol turns one such statement into a smaller one of
the same shape, and the proof is the transcript of repeating that until the
statement is small enough to answer in the clear.

\begin{enumerate}
  \item \textbf{Commit.} The prover commits to each $\mathbf{s}_i$ with an
        Ajtai commitment $\mathbf{t}_i=\mathbf{A}\mathbf{s}_i$
        (Figure~\ref{fig:ajtai-commit}), then commits again to the vector of
        all the $\mathbf{t}_i$, split into short digits, so that what is sent
        is one short \emph{outer} commitment rather than $r$ inner ones. The
        same is done for the inner products
        $g_{ij}=\langle\mathbf{s}_i,\mathbf{s}_j\rangle$ that the quadratic
        constraints will need.
  \item \textbf{Prove the norm by projection.} The verifier sends random
        matrices $\Pi_i$ with entries in $\{-1,0,1\}$ and the prover returns
        the projection $\mathbf{p}=\sum_i\Pi_i\mathbf{s}_i$, a vector of $256$
        integers regardless of $r$ and $n$. By the Johnson--Lindenstrauss
        lemma, $\|\mathbf{p}\|$ is, with overwhelming probability, within a
        constant factor of $\sqrt{\smash[b]{\textstyle\sum_i\|\mathbf{s}_i\|^2}}$,
        so a short $\mathbf{p}$ certifies the norm bound at a cost independent
        of the witness size. Verifying that $\mathbf{p}$ really is the
        projection is itself a linear constraint on the $\mathbf{s}_i$, and is
        folded into the next step.
  \item \textbf{Aggregate and amortize.} The verifier sends random weights, and
        all the constraints are combined into a few. It then sends short
        polynomial challenges $c_1,\dots,c_r$, and the prover answers with the
        single vector $\mathbf{z}=\sum_i c_i\mathbf{s}_i$. The verifier checks
        $\mathbf{A}\mathbf{z}=\sum_i c_i\mathbf{t}_i$, which the homomorphism
        makes possible without seeing any $\mathbf{s}_i$, and checks the
        combined constraint against $\mathbf{z}$ and the committed $g_{ij}$,
        using $\langle\mathbf{z},\mathbf{z}\rangle=\sum_{i,j}c_ic_j\,g_{ij}$.
  \item \textbf{Recurse.} What the verifier now holds is a claim that
        $\mathbf{z}$, together with the openings of the outer commitments, is
        short and satisfies a dot-product constraint system --- the same shape
        as the input, with the $r$ witness vectors replaced by one (plus the
        digit vectors of the commitments), and with the norm bound grown by
        the size of the challenges. Instead of sending $\mathbf{z}$, the prover
        treats this new statement as the input to another round.
\end{enumerate}

Figure~\ref{fig:labrador-round} writes one round in the convention of the
rest of this section.

\begin{figure}[H]
\begin{protocol}
\textbf{Public information:} $\mathbf{A}\in R_q^{\kappa\times n}$ and
$\mathbf{B},\mathbf{B}'$ for the outer commitments; the constraint system
$\{(a^{(k)}_{ij},\boldsymbol{\varphi}^{(k)}_i,b^{(k)})\}_k$, where
$a^{(k)}_{ij}\in R_q$ weights $\langle\mathbf{s}_i,\mathbf{s}_j\rangle$,
$\boldsymbol{\varphi}^{(k)}_i\in R_q^{n}$ is the linear coefficient of
$\mathbf{s}_i$, and $b^{(k)}\in R_q$ is the constant; norm bound
$\beta$; a challenge set $\mathcal{C}\subset R$ of short polynomials.\\
\textbf{Private information:} $\mathbf{s}_1,\dots,\mathbf{s}_r\in R^{n}$ with
$\sum_i\|\mathbf{s}_i\|^2\le\beta^2$, satisfying every constraint
$\sum_{i,j}a^{(k)}_{ij}\langle\mathbf{s}_i,\mathbf{s}_j\rangle
+\sum_i\langle\boldsymbol{\varphi}^{(k)}_i,\mathbf{s}_i\rangle+b^{(k)}=0$.
\protorule
\begin{tabular}{@{}p{0.34\linewidth}@{}p{0.26\linewidth}@{}p{0.40\linewidth}@{}}
\Prover & & \Verifier \\[0.6em]
$\mathbf{t}_i:=\mathbf{A}\mathbf{s}_i$ for all $i$; & & \\
$\mathbf{u}_1:=\mathbf{B}\,G^{-1}(\mathbf{t}_1\|\cdots\|\mathbf{t}_r)$ & & \\[0.3em]
 & \centering$\msgto{\mathbf{u}_1}$ & \\[0.3em]
 & \centering$\msgfrom{\Pi_1,\dots,\Pi_r}$ & $\Pi_i\leftarrow\{-1,0,1\}^{256\times nd}$ \\[0.3em]
$\mathbf{p}:=\sum_i\Pi_i\,\mathbf{s}_i\in\mathbb{Z}^{256}$ & & \\[0.3em]
 & \centering$\msgto{\mathbf{p}}$ & \\[0.3em]
 & \centering$\msgfrom{\text{weights}}$ & fold all constraints, and \\
 & & ``$\mathbf{p}$ is the projection'', into one \\[0.3em]
$g_{ij}:=\langle\mathbf{s}_i,\mathbf{s}_j\rangle$,\ $h_{ij}$ from the folded linear parts; & & \\
$\mathbf{u}_2:=\mathbf{B}'\,G^{-1}(g\|h)$ & & \\[0.3em]
 & \centering$\msgto{\mathbf{u}_2}$ & \\[0.3em]
 & \centering$\msgfrom{c_1,\dots,c_r}$ & $c_i\leftarrow\mathcal{C}$ \\[0.3em]
$\mathbf{z}:=\sum_i c_i\,\mathbf{s}_i$ & & \\[0.3em]
 & \centering$\msgto{\mathbf{z},\ G^{-1}(\mathbf{t}),\ G^{-1}(g\|h)}$ & \\[0.3em]
 & & Accept iff \\
 & & \quad 1. $\|\mathbf{p}\|\le\sqrt{128}\,\beta$, and $\mathbf{z}$ and both digit vectors are short \\
 & & \quad 2. $\mathbf{B}\,G^{-1}(\mathbf{t})=\mathbf{u}_1$,\ $\mathbf{B}'\,G^{-1}(g\|h)=\mathbf{u}_2$ \\
 & & \quad 3. $\mathbf{A}\mathbf{z}=\sum_i c_i\,\mathbf{t}_i$ \\
 & & \quad 4. $\langle\mathbf{z},\mathbf{z}\rangle=\sum_{i,j}c_ic_j\,g_{ij}$ \\
 & & \quad 5. $\bigl\langle\sum_i c_i\boldsymbol{\varphi}_i,\ \mathbf{z}\bigr\rangle=\sum_{i,j}c_ic_j\,h_{ij}$ \\
 & & \quad 6. the folded constraint is zero on $(g,h)$ \\
\end{tabular}
\end{protocol}
\caption{One round of LaBRADOR, in teaching form. Read as a stand-alone
protocol, the last message opens everything and the verifier's checks are as
listed. In the recursive protocol the last message is \emph{not} sent: checks
1--6 are a statement about the short vectors $\mathbf{z}$, $G^{-1}(\mathbf{t})$
and $G^{-1}(g\|h)$ that is again a dot-product constraint system, and it becomes
the input to the next round. Check 1 is the Johnson--Lindenstrauss step: $256$
integers stand in for the norm of the whole witness. Check 5 is where the
verifier's linear cost enters, since $\sum_ic_i\boldsymbol{\varphi}_i$ ranges
over the original public statement.}
\label{fig:labrador-round}
\end{figure}

Each round shrinks the witness by roughly the factor $r$ that amortization
buys, and the total transcript is a handful of commitments and one
$256$-entry projection per round, plus the final small witness sent in the
clear. That is where the size comes from: about $58$\,kB for an R1CS instance
of $2^{20}$ constraints at the 128-bit level, essentially independent of the
statement size beyond a few rounds. And it is where the verifier's cost comes
from: to check the combined constraint the verifier must form
$\sum_i c_i\boldsymbol{\varphi}_i$ over the \emph{original} public statement,
so its work is linear in the statement even though the proof is not. Greyhound
(below) is the construction that supplies a sublinear verifier, and it uses
LaBRADOR only to compress its transcript.

The outline above compresses two papers into a page. A lecture-length treatment
of the same path --- Ajtai commitments and Lyubashevsky-style proofs, then
LaBRADOR, then Greyhound, then the LaBRADOR and Greyhound proof-of-concept
implementations --- with slides, recordings and participants' notes, is the
lattice course of the ZKP Academy~\cite{zkpacademylattice}, and it is the
recommended companion to this section for a reader who wants the protocols in
full rather than in outline.

Extraction from such a proof carries a cost. It does not hand back the witness
that was committed. What it yields is a
$\bar{\mathbf{s}}$ and a short polynomial $c$ with
\[
  \mathbf{A}\bar{\mathbf{s}} = c\,\mathbf{t} \bmod q,
  \qquad \|\bar{\mathbf{s}}\| \ \text{larger than}\ \|\mathbf{s}\| ,
\]
which is a different equation: the norm has grown, and an unknown multiplicand
has appeared. That gap is called \textbf{slack}, and it is not a proof artefact
to be tidied away later.

Two consequences show why. Practically, a scheme downstream of such a proof
inherits the $c$. In lattice verifiable encryption this was for years the reason
the decryptor had to \emph{guess} $c$ in order to decrypt, making expected
decryption time track the adversary's running time --- so the scheme could only
be deployed where an attacker's time could be bounded~\cite{lnp2022}. And in
parameters, slack compounds: a larger extracted solution demands a larger modulus
for SIS security, which demands a larger dimension for LWE security, which
enlarges everything again. It is the same wall that aggregating Falcon signatures
hit, where LaBRADOR's norm check proved ``both approximate and with respect to
the entire witness,'' with the slack growing in the number of signatures, forcing
the first iteration to be replaced with an exact proof of
smallness~\cite{falconagg2024}. Much of the work in this area over the past
decade concerns reducing or removing slack.

\subsubsection{From a commitment to a polynomial commitment}
Figure~\ref{fig:ajtai-commit} commits to a \emph{vector}. A proof system needs
to commit to a \textbf{polynomial} $f$ and later convince a verifier that
$f(x)=y$ at a point $x$ chosen after the commitment --- a \textbf{polynomial
commitment scheme}, or PCS. Two requirements: it is \emph{binding}, so no prover
can answer two different values at the same point; and it is \emph{succinct}, so
the commitment and the evaluation proof are far smaller than $f$ itself.

\paragraph{The polynomial commitment as the unit of migration.}
Most modern argument systems --- STARKs, PLONK and their descendants
--- factor into two independent pieces:
\begin{enumerate}
  \item a \textbf{polynomial interactive oracle proof}, in which the prover hands
        the verifier idealized oracles for some polynomials and the verifier
        queries them. This half is \emph{information-theoretic}: it assumes
        nothing, and a quantum adversary gains nothing against it.
  \item a \textbf{compiler} that replaces each oracle with a polynomial
        commitment, and each interaction with Fiat--Shamir.
\end{enumerate}
Every computational assumption in the system enters through the second step.
That is why ``is this proof system post-quantum?'' is nearly always a question
about its PCS, and why the Ziren repair earlier in this section is a primitive
swap rather than a redesign: the protocol above the commitment does not know or
care what the commitment is made of.

\begin{table}[H]
\centering
\small
\begin{tabular}{@{}p{2.3cm}p{3.5cm}p{3.4cm}p{2.2cm}p{2.4cm}@{}}
\toprule
PCS & Commitment is & Rests on & Setup & Survives Shor? \\
\midrule
KZG & one group element & pairings; discrete log & structured, trusted & \textbf{no} \\
FRI / WHIR & a Merkle root & random oracle; collision resistance & transparent & yes \\
Ajtai & $\mathbf{t}=\mathbf{A}\mathbf{s}$, $\kappa$ ring elements & Module-SIS (+ random oracle once compiled) & transparent & yes \\
\bottomrule
\end{tabular}
\caption{The three polynomial commitment families, which is the row that decides
whether an argument system is post-quantum. Ethereum's stated plan to ``move KZG
commitments to STARK- or lattice-based alternatives'' (Section~\ref{sec:eth-two-front})
is a move from the first row to one of the others; the protocol layered on top is
largely unaffected.}
\label{tab:pcs-families}
\end{table}

\paragraph{The square-root decomposition.}
For the lattice row, the bridge from Figure~\ref{fig:ajtai-commit} to a PCS is
a reshaping. Write
$f(X)=\sum_{i<N} f_i X^{i}$ with $N=n\cdot m$, arrange the coefficients into an
$n\times m$ matrix $M$ with $M_{i,j}=f_{i+jn}$, and the evaluation factors:
\[
  f(x)\;=\;\sum_{j<m} x^{jn} \sum_{i<n} f_{i+jn}\,x^{i}
       \;=\;\mathbf{a}(x)^{\!\top} M\, \mathbf{b}(x),
  \qquad
  \begin{aligned}
    \mathbf{a}(x)&=(1,x,\dots,x^{n-1}),\\
    \mathbf{b}(x)&=(1,x^{n},\dots,x^{(m-1)n}).
  \end{aligned}
\]
A quadratic form, and a commitment that is linearly homomorphic can be pushed
through one side of it.

\begin{figure}[H]
\begin{protocol}
\textbf{Public information:} $\mathbf{A} \in R_q^{\kappa\times n}$, the column
commitments $\mathbf{t}_j = \mathbf{A}M_{\cdot,j}$ for $j<m$, the evaluation
point $x$, the claimed value $y$, a norm bound $B$.\\
\textbf{Private information:} $M \in R^{n\times m}$, the coefficients of $f$,
such that
\[
  \mathbf{a}(x)^{\!\top} M\, \mathbf{b}(x) = y .
\]
\protorule
\begin{tabular}{@{}p{0.30\linewidth}@{}p{0.28\linewidth}@{}p{0.38\linewidth}@{}}
\Prover & & \Verifier \\[0.9em]
$\mathbf{w} := M\,\mathbf{b}(x)$ & & \\[0.6em]
 & \centering$\msgto{\mathbf{w}}$ & \\[0.6em]
 & & Accept iff \\
 & & \quad 1. $\|\mathbf{w}\| \le B$ \\
 & & \quad 2. $\langle \mathbf{a}(x), \mathbf{w}\rangle = y$ \\
 & & \quad 3. $\mathbf{A}\mathbf{w} = \sum_{j<m} b_j(x)\,\mathbf{t}_j$ \\
\end{tabular}
\end{protocol}
\caption{The square-root evaluation argument, in the shape Greyhound
instantiates. Check 3 is the homomorphism doing the work: the verifier never
sees the columns, but
$\sum_j b_j \mathbf{A}M_{\cdot,j} = \mathbf{A}(M\mathbf{b}) = \mathbf{A}\mathbf{w}$,
so it can derive the commitment the response \emph{should} have and compare.
Check 1 is not a formality, for the reason given below. As drawn this is one round
and it is not yet sound --- see the caveat below.}
\label{fig:sqrt-pcs}
\end{figure}

Both the transcript and the verifier are $O(\sqrt{N})$: the prover sends one
column's worth of data instead of the whole matrix, and the verifier does
$O(\sqrt{N})$ ring operations instead of reading $f$.

\paragraph{Why the real protocol has three rounds.}
The figure is the idea, and one step short of a protocol. The entries of
$\mathbf{b}(x)$ are powers of $x$ and are not small, so $\mathbf{w}=M\mathbf{b}(x)$
is not short either and check 1 cannot be satisfied at any useful bound. A real
instantiation therefore inserts a round: the verifier sends a \emph{short}
random challenge, the prover answers with a combination weighted by it, and the
response stays inside the norm budget. That extra round is the difference between
the picture and Greyhound; the difficulty lies in the norm bounds rather than
in the algebra. Figure~\ref{fig:greyhound-eval} gives the protocol as Greyhound
runs it. It differs from the sketch in two further respects: what is public,
and which side of the quadratic form the challenge lands on.

\paragraph{Inner and outer commitments.}
Figure~\ref{fig:sqrt-pcs} lists the column commitments $\mathbf{t}_j$ as public.
In Greyhound they are not; only a commitment \emph{to} them is. The reason is
size. Each $\mathbf{t}_j=\mathbf{A}\mathbf{f}_j$ is a binding Ajtai commitment
to a short column, but there are $m=O(\sqrt{N})$ of them, so a commitment to
$f$ consisting of all the $\mathbf{t}_j$ would be $O(\sqrt{N})$ long. Greyhound
therefore commits twice: the \textbf{inner} commitments
$\mathbf{t}_j=\mathbf{A}\mathbf{f}_j$, and an \textbf{outer} commitment
\[
  \mathbf{u}\;=\;\mathbf{B}\,G^{-1}(\mathbf{t}_1\|\cdots\|\mathbf{t}_m)
\]
to the concatenation of all of them. The gadget inverse $G^{-1}$ is not
optional: the $\mathbf{t}_j$ are uniform modulo $q$, and an Ajtai commitment
binds only short inputs, so they are first split into short digits and the
digit vector is what $\mathbf{B}$ commits. The polynomial commitment is
$\mathbf{u}$, a handful of ring elements whatever $N$ is; the $\mathbf{t}_j$
are revealed inside the evaluation proof and checked against it. This is the
same two-level pattern as step 1 of the LaBRADOR outline in
Section~\ref{sec:labrador-outline}.

\paragraph{Where the challenge lands.}
The response that must be short is a combination of \emph{columns} weighted
by a short challenge, $\mathbf{z}=F\mathbf{c}$; the partial evaluations
$w_j=\langle\mathbf{a}(x),\mathbf{f}_j\rangle$ are sent in the clear and need
not be short at all, because $\mathbf{z}$ pins them.

\begin{figure}[H]
\begin{protocol}
\textbf{Public information:} $\mathbf{A}\in R_q^{\kappa\times n}$,
$\mathbf{B}$, the outer commitment $\mathbf{u}$, the evaluation point $x$, the
claimed value $y$, a norm bound $B_z$, a challenge set $\mathcal{C}$.\\
\textbf{Private information:} short columns
$F=[\mathbf{f}_1|\cdots|\mathbf{f}_m]\in R^{n\times m}$ with
$\mathbf{a}(x)^{\!\top}F\,\mathbf{b}(x)=y$, and
$\mathbf{u}=\mathbf{B}\,G^{-1}(\mathbf{t})$ for
$\mathbf{t}_j=\mathbf{A}\mathbf{f}_j$.
\protorule
\begin{tabular}{@{}p{0.34\linewidth}@{}p{0.26\linewidth}@{}p{0.40\linewidth}@{}}
\Prover & & \Verifier \\[0.6em]
$\mathbf{d}:=G^{-1}(\mathbf{t}_1\|\cdots\|\mathbf{t}_m)$; & & \\
$w_j:=\langle\mathbf{a}(x),\mathbf{f}_j\rangle$ for $j<m$ & & \\[0.3em]
 & \centering$\msgto{\mathbf{d},\ \mathbf{w}}$ & \\[0.3em]
 & \centering$\msgfrom{\mathbf{c}}$ & $\mathbf{c}=(c_1,\dots,c_m)\leftarrow\mathcal{C}^m$ \\[0.3em]
$\mathbf{z}:=F\mathbf{c}=\sum_j c_j\,\mathbf{f}_j$ & & \\[0.3em]
 & \centering$\msgto{\mathbf{z}}$ & \\[0.3em]
 & & Accept iff \\
 & & \quad 1. $\mathbf{d}$ is a valid digit vector and $\|\mathbf{z}\|\le B_z$ \\
 & & \quad 2. $\mathbf{B}\mathbf{d}=\mathbf{u}$; set $\mathbf{t}:=G\mathbf{d}$ \\
 & & \quad 3. $\langle\mathbf{w},\mathbf{b}(x)\rangle=y$ \\
 & & \quad 4. $\mathbf{A}\mathbf{z}=\sum_j c_j\,\mathbf{t}_j$ \\
 & & \quad 5. $\langle\mathbf{a}(x),\mathbf{z}\rangle=\sum_j c_j\,w_j$ \\
\end{tabular}
\end{protocol}
\caption{Greyhound's evaluation protocol, in teaching form. Check 2 opens the
outer commitment and recovers the inner ones; check 4 is the homomorphism of
Figure~\ref{fig:sqrt-pcs} applied to a \emph{short} challenge combination, so
the norm bound in check 1 can be met; check 5 ties the claimed partial
evaluations to that same combination, and check 3 finishes the evaluation from
them. Prover message and verifier work are both $O(m+n)=O(\sqrt{N})$. The
composition with LaBRADOR proves that a transcript satisfying checks 1--5
exists, instead of sending $\mathbf{d}$, $\mathbf{w}$ and $\mathbf{z}$.}
\label{fig:greyhound-eval}
\end{figure}

\subsubsection{SageMath experiment: the commitment and the evaluation protocol}
The two figures are short enough to run. The first listing is
Figure~\ref{fig:ajtai-commit} in a toy ring, $R_q=\mathbb{Z}_{12289}[x]/(x^{16}+1)$
with $\kappa=2$ and $n=8$: a uniformly random $\mathbf{A}$, a ternary witness,
the commitment $\mathbf{t}=\mathbf{A}\mathbf{s}$, and the verifier's two checks.
The parameters are far too small for security; they are chosen so that every
object can be printed.

\begin{lstlisting}[language=Python]
# The Ajtai commitment in a toy ring: t = A s with s short.
q, d, kappa, n = 12289, 16, 2, 8
R.<x> = PolynomialRing(GF(q)); Rq.<X> = R.quotient(x^d + 1)

def centered(v):        # coefficients of a ring element in (-q/2, q/2]
    return [c - q if c > q // 2 else c for c in map(Integer, v.lift().list() + [0]*d)][:d]
def norm2(vec):         # squared Euclidean norm of a vector of ring elements
    return sum(c^2 for v in vec for c in centered(v))
def short(k, bound=1):  # k ring elements with coefficients in [-bound, bound]
    return vector(Rq, [Rq([randint(-bound, bound) for _ in range(d)]) for _ in range(k)])

A = random_matrix(Rq, kappa, n)      # public randomness, no trapdoor
s = short(n)                          # the witness
t = A * s                             # the commitment: kappa ring elements
beta2 = n * d                         # every coefficient at most 1 in size

def open_ok(t, s):                    # the verifier of Figure 28.2
    return A * s == t and norm2(s) <= beta2

print("opens:", open_ok(t, s))
s_long = short(n, bound=q // 4)       # long vectors fail the norm check, whatever A s is
print("long vector passes the norm check:", norm2(s_long) <= beta2)
\end{lstlisting}

The second listing continues from the first and runs
Figure~\ref{fig:greyhound-eval}: an $n\times m$ matrix $F$ of short ring
elements, the $m$ inner commitments, the gadget decomposition into base-16
digits and the outer commitment $\mathbf{u}$, then the three rounds and the
five checks. \texttt{digits} is $G^{-1}$ and \texttt{undigits} is $G$; the
verifier never sees $F$ or the $\mathbf{t}_j$ except through
$\mathbf{u}$ and the digits it is handed. Tampering with one partial
evaluation makes check~5 fail, which is the consistency that $\mathbf{z}$
enforces.

\begin{lstlisting}[language=Python]
# The square-root PCS: inner and outer commitments, then the
# three-round evaluation protocol.  Continues from the Ajtai listing.
m = 4                                  # F is n x m: N = n*m short ring elements
ell, base = 4, 16                      # 16^4 > q: four digits per coefficient

def digits(vec):                       # G^{-1}: split each coefficient in base 16
    out = []
    for v in vec:
        cs = [Integer(c) for c in v.lift().list() + [0]*d][:d]
        for j in range(ell):
            out.append(Rq([(c >> (4*j)) & 15 for c in cs]))
    return vector(Rq, out)
def undigits(dv):                      # G: the inverse map
    return vector(Rq, [sum(dv[i*ell + j] * base^j for j in range(ell))
                       for i in range(len(dv) // ell)])

B = random_matrix(Rq, kappa, kappa * m * ell)   # for the outer commitment
F = matrix(Rq, n, m, lambda i, j: short(1)[0])  # the committed coefficients
T = [A * F.column(j) for j in range(m)]         # inner commitments t_j
tvec = vector(Rq, sum((list(t) for t in T), []))
u = B * digits(tvec)                             # outer commitment: the PCS commitment

xpt = Rq.random_element()                        # evaluation point
a = vector(Rq, [xpt^i for i in range(n)]); b = vector(Rq, [xpt^(n*j) for j in range(m)])
y = a * F * b                                    # the claimed value f(x)

# Round 1: prover sends the digits and the partial evaluations.
dvec = digits(tvec); w = vector(Rq, [a * F.column(j) for j in range(m)])
# Round 2: verifier sends a short challenge.  Round 3: prover answers with z.
c = short(m); z = F * c

def verify(dvec, w, c, z):                       # the five checks of the evaluation figure
    digit_ok = all(0 <= Integer(v) < base for dd in dvec for v in dd.lift().list())
    tv = undigits(dvec); tj = [tv[i*kappa:(i+1)*kappa] for i in range(m)]
    return (digit_ok and norm2(z) <= n * d * (m * d)^2   # 1. digits and shortness
            and B * dvec == u                    # 2. outer commitment opens
            and w * b == y                       # 3. the evaluation
            and A * z == sum(c[j] * tj[j] for j in range(m))   # 4. inner homomorphism
            and a * z == sum(c[j] * w[j] for j in range(m)))   # 5. w is consistent with z

print("honest transcript verifies:", verify(dvec, w, c, z))
w_bad = vector(Rq, list(w)); w_bad[0] += 1
print("with a tampered partial evaluation:", verify(dvec, w_bad, c, z))
\end{lstlisting}

What the experiment does not do is compress: the prover sends $\mathbf{d}$,
$\mathbf{w}$ and $\mathbf{z}$ in the clear, $O(\sqrt{N})$ ring elements, and
the verifier does $O(\sqrt{N})$ ring multiplications. Greyhound's remaining
step is to hand the five checks, as a dot-product constraint system in
$\mathbf{d}$ and $\mathbf{z}$, to LaBRADOR (Section~\ref{sec:labrador-outline}),
which is what shrinks the transcript to tens of kilobytes.

\paragraph{The role of the norm check.}
Without the norm check there is no commitment at all. The map $\mathbf{s}\mapsto\mathbf{A}\mathbf{s}$ from $R_q^{n}$ to
$R_q^{\kappa}$ with $n\gg\kappa$ is wildly non-injective --- every $\mathbf{t}$
has on the order of $q^{d(n-\kappa)}$ preimages, and Gaussian elimination finds
one in polynomial time. Anyone can open an Ajtai commitment to anything they
like, provided they are allowed to answer with a large vector.

What is hard is finding a \emph{short} preimage, and finding \emph{two} short
preimages of the same $\mathbf{t}$ is exactly the Module-SIS problem of the
binding argument above. So the norm bound is not a side condition on the scheme;
it \emph{is} the scheme. Every message in every round carries one, every
operation that combines messages has to account for the growth, and a protocol
that loses track of a bound anywhere has lost binding everywhere. This is why
lattice arguments read like an exercise in bookkeeping, and why the extraction
slack of the previous subsection is the natural failure mode rather than a
surprise.

\paragraph{Comparison with FRI.}
The mechanism differs from FRI's more than their shared post-quantum label
suggests. FRI never proves an algebraic identity: it establishes that a committed word is \emph{close} to a low-degree
codeword, by folding the word and opening a few positions, and its cost is
therefore counted in queries and Merkle paths --- which is precisely why a
verifier written in Bitcoin Script spends nearly all of its budget on hashing.
The lattice PCS proves the identity directly, and its cost is counted in ring
multiplications and in norm growth. The two are different mechanisms with
different cost centres; replacing one with the other changes the system's cost
profile, not a parameter.

\paragraph{Composition.}
Greyhound's contribution is the three-round protocol of
Figure~\ref{fig:greyhound-eval} with the norm control that lattices require,
and then a \emph{composition}: checks 1--5 of that figure are linear and
quadratic relations in the short vectors $\mathbf{d}$ and $\mathbf{z}$, which is
a dot-product constraint system in the sense of
Section~\ref{sec:labrador-outline}, so LaBRADOR can prove that a satisfying
transcript exists and compress the $O(\sqrt{N})$ transcript to
polylogarithmic size~\cite{greyhound2024}.

The composition affects only the proof size. The sublinear verifier is already
present in the base protocol; LaBRADOR shrinks the proof and leaves the
verifier's cost unchanged, at $O(\sqrt{N})$.

\subsubsection{What the verifier actually computes}
Verification is often summarized as ring multiplication by way of the NTT; the
implementation is less direct. A proof modulus is
chosen to make the \emph{security} argument work, not to be NTT-friendly, and
Greyhound's is not: its implementation states plainly that NTT-based
multiplication cannot be used directly for $R_q$, and instead multiplies modulo
several small primes of 12 to 14 bits that \emph{do} split fully in
$\mathbb{Z}[X]/(X^{64}+1)$, recombining by explicit CRT~\cite{greyhound2024}.

The transform of Chapter~\ref{ch:ntt} still runs, but several times per
multiplication, followed by a CRT reconstruction. For a verifier that pays per
operation, this constant factor is part of the design.

\paragraph{Proof size against verification cost.}
Lattice arguments produce the smallest post-quantum proofs by a wide margin: LaBRADOR proves an
R1CS with $2^{20}$ constraints in \textbf{58\,kB} at the 128-bit level, against
megabytes for the hash-based family. The cost has moved rather than disappeared,
and it has moved to the verifier --- LaBRADOR's is \emph{linear} in the witness,
and Greyhound measures \textbf{2.80\,s} to verify a degree-$2^{30}$ evaluation,
which its own authors put at roughly twice Ligero's~\cite{greyhound2024}. A
verifier living inside a contract or a script is buying bytes with time, and
should confirm it wants that trade.

Published figures in this line also need reading carefully, because no two of
them measure the same object: Greyhound's headline 53\,kB is an \emph{evaluation}
proof rather than a proof of a computation, and the much-quoted 73\,kB from
SALSAA is one folding step over linear relations, against 979\,kB for the same
system's general-purpose argument~\cite{salsaa2025}.

\paragraph{Signature aggregation.}
No standardized post-quantum signature aggregates
(Section~\ref{sec:symmetric-plumbing}); one research result exists. Falcon
signatures can be aggregated by proving their
verification inside LaBRADOR: 500 signatures in about 73\,kB, against roughly
3.6\,MB for the best lattice aggregate-signature scheme
otherwise~\cite{falconagg2024}. Two qualifications apply. The parameters are at 121-bit security rather than 128, and every
size is an analytic estimate --- there is no prover or verifier implementation
to measure.

The mechanism also shows why aggregation is not free even when it works. Beyond
the exact-smallness repair above, the proof cannot run over Falcon's own modulus
at all: the projection the argument uses needs room that $q=12289$ does not have,
so the statement is lifted to a proof modulus of \emph{61 bits} for FN-DSA-512
and 63 for FN-DSA-1024~\cite{falconagg2024}. Aggregating signatures means proving
statements about them in a different arithmetic from the one they were signed
in.

A third qualification concerns generality. The result \textbf{does not
transfer to ML-DSA}, for a structural reason: its soundness proof
handles a signing oracle only for \emph{hash-then-sign} schemes, the GPV
structure of Chapter~\ref{ch:fndsa}, and ML-DSA's Fiat--Shamir-with-aborts
structure (Chapter~\ref{ch:mldsa}) is exactly what the proof cannot
accommodate. The one existing aggregation result therefore applies to FN-DSA
and not to ML-DSA, which adds a consideration beyond size to the comparison of
Section~\ref{sec:bitcoin-scheme-choice}.

\paragraph{Assumptions within the lattice family.}
The family label determines as little here as it did in
Section~\ref{sec:symmetric-plumbing}. In 2024 a quantum polynomial-time
algorithm was shown to sample well-distributed LWE
instances while provably \emph{not} knowing the secret, assuming only that LWE
is hard~\cite{obliviouslwe2024}. It invalidated the security analyses of six
published lattice SNARKs at a stroke.

The scope of the result is instructive. The six affected schemes are
standard-model, designated-verifier constructions resting on a
\emph{knowledge-of-encoding} assumption --- an assumption that the sampler
refutes. The LaBRADOR line is untouched, because it assumes Module-SIS in the
random oracle model and never assumes that an adversary producing an encoding
must know what it encoded. The two lines share a primitive family but rest on
different assumptions, with different outcomes; the paper states that it
invalidated security analyses rather than broke the constructions. At every
level of the stack, the assumption rather than the family label is what
determines the security.

\section{Timelines and honest uncertainty}
No quantum computer capable of breaking secp256k1 or BLS12-381 exists as of mid-2026.
A 2026 arXiv preprint, not a consensus survey, gives one set of model-dependent
estimates for cryptographically relevant quantum computing~\cite{quantumhorizon2026}.
A 2025 Google analysis of RSA-2048 illustrates a downward resource-estimation trend,
but is RSA-specific and not a direct estimate for Bitcoin or Ethereum curves~\cite{gidney2025}.
Expert elicitation~\cite{mosca2025} and national guidance~\cite{nistir8547,ncscpqc2025,cnsa2}
place the migration deadline for conventional systems in the early-to-mid 2030s,
but those targets were set for enterprise IT, where an administrator can deploy a
change; they were not calibrated for a system whose users must each act
individually.

The honest position is that $z$ in Mosca's inequality is unknown and contested.
The useful response is to stop arguing about $z$ and measure $y$ instead --- the
migration time is a property of the system, not of the adversary, and unlike $z$
it can be computed.

\subsection{Migration is bounded by blockspace}
\label{sec:migration-throughput}

Here is the constraint that the roadmap debates tend to skip. Migrating a coin is
not a database update; it is a \emph{transaction}, and transactions compete for a
fixed and quite small amount of blockspace. So there is a floor on how fast any
migration can possibly proceed, and it depends only on arithmetic --- the number
of outputs to move, the size of a post-quantum spend, and the throughput of the
chain.

Write $N$ for the number of unspent outputs to migrate, $w$ for the weight a
single migration transaction consumes, $W$ for the weight available per block,
and $T$ for the block interval. If migration were given \emph{every unit} of
every block, it would still take
\[
  T_{\min} \;=\; T \cdot \left\lceil \frac{N\,w}{W} \right\rceil ,
\]
and at a share $\sigma$ of blockspace the figure scales by $1/\sigma$.

Two features of this bound deserve emphasis. It is a \emph{lower} bound assuming
all of the chain's capacity goes to migration and nothing else --- no payments,
no ordinary economic activity --- which no live chain would tolerate. And it is
driven by $w$, which for a post-quantum spend is dominated by the signature and
public key, exactly the quantities that are one to two orders of magnitude larger
than what they replace. The sizes tabulated throughout this book are not merely a
bandwidth annoyance: they set the clock on the migration itself.

A published analysis of this question for Bitcoin reports a required downtime on
the order of a few months of exclusive blockspace, under its own assumptions
about which outputs must move and what a migration spend
costs~\cite{pont2024}. The exercise below builds a cruder model and gets a
larger answer; the discrepancy is the instructive part, and reconciling the two
is the point.

\begin{remark}[Why this reframes the whole chapter]
Every other quantity in this chapter is contested: when Q-Day arrives, whether a
freeze is legitimate, which signature scheme wins. The throughput bound is not.
It follows from block size, signature size, and arithmetic, and it says that even
a perfectly coordinated ecosystem with universal agreement and instant wallet
support could not complete the migration quickly. That is the strongest available
argument for starting early, and it does not require winning any argument about
quantum hardware.
\end{remark}

\subsection*{Exercise: compute the bound}

The sensitivity of the formula to $w$ is best seen by evaluating it. Using the tools of Chapter~\ref{ch:sagemath}:

\begin{lstlisting}[language=Python]
# Lower bound on migration time: every weight unit of every block, nothing else.
BLOCK_WEIGHT   = 4_000_000   # weight units available per Bitcoin block
BLOCKS_PER_DAY = 144         # one block per ten minutes

def migration_days(n_outputs, witness_bytes, base_bytes=80, share=1.0):
    # SegWit discount: witness bytes cost 1 weight unit, the rest cost 4.
    weight_per_tx = witness_bytes + 4 * base_bytes
    per_block = BLOCK_WEIGHT // weight_per_tx
    blocks = -(-n_outputs // per_block)          # ceiling division
    return blocks / (BLOCKS_PER_DAY * share)

SCHEMES = [("ed25519 (baseline)",  64 + 32),      # signature + public key
           ("FN-DSA-512",         666 +  897),
           ("FN-DSA-1024",       1280 + 1793),
           ("ML-DSA-65",         3309 + 1952),
           ("SLH-DSA-128s",      7856 + 32)]

for label, n in [("every UTXO (~190M)", 190_000_000),
                 ("exposed third (~63M)", 63_000_000)]:
    print(label)
    for name, size in SCHEMES:
        print(f"  {name:20s} {migration_days(n, size):8.0f} days exclusive"
              f"   {migration_days(n, size, share=0.10):8.0f} days at 10%")
\end{lstlisting}

The model gives roughly 1800 days to move every output at ML-DSA-65 sizes with
the whole chain to itself, and about 600 days for the exposed third --- against
137 and 46 days respectively for the classical baseline. At a tenth of
blockspace, which is already generous for a chain that must keep processing
payments, the exposed third alone runs to some sixteen years.

Four things to take from that. The baseline rows show that migration is slow
even with \emph{classical} signatures, so post-quantum sizes are an aggravating
factor rather than the sole cause. The SLH-DSA rows price the most conservative
assumption in schedule rather than in bytes --- a trade-off the size tables of
Chapter~\ref{ch:applications} state but do not cost out. The FN-DSA rows are the
reason Section~\ref{sec:bitcoin-scheme-choice} cares about the combined
footprint: at Falcon-512 sizes the schedule is roughly a third of the ML-DSA-65
figure and a fifth of the SLH-DSA one, so the scheme choice moves the migration
clock by years, not percentages. And the gap between
these figures and the published estimate is itself the lesson: this model
migrates one output per transaction, at full witness size, for the entire
exposed set. Every one of those is pessimistic, and the published analysis relaxes
them.

So relax them, and see which matters. Vary the migrating set; batch many outputs
into one transaction; apply an aggregation or a proof-based authorization in
place of a per-output signature. Batching is the interesting lever, because $N$
and $W$ are given to the designer while $w$ per migrated output is the one term a
protocol can actually negotiate --- which is precisely what the output-type and
recovery-path proposals in this chapter are trying to buy.

\begin{remark}[Do not quote this figure]
The numbers above depend on the UTXO count, the fraction deemed exposed, the
witness discount, batching, and the scheme --- and Pitfall 3 of this chapter
applies to them as much as to gas figures. The claim that survives the
uncertainty is the order of magnitude: \emph{years, not weeks}, under any
assumption set that keeps the chain usable.
\end{remark}

\section{Pitfalls}

\paragraph{Pitfall 1: assuming ``my address is a hash, so I am safe.''}
True only until the first spend. Any account that has transacted has exposed its public key; on Ethereum that is essentially every active account, and on Bitcoin it includes every reused address.

\paragraph{Pitfall 2: treating draft proposals as decided.}
BIP-360, BIP-361, and EIP-8141 are drafts; the freeze question in particular is contested. Reporting them as Bitcoin's or Ethereum's ``plan'' overstates consensus.

\paragraph{Pitfall 3: quoting a single gas figure.}
On-chain verification costs move between implementations, versions, and precompile assumptions. A number without a source and date is not meaningful.

\paragraph{Pitfall 4: assuming hybrid signatures are the safe default.}
They are the conservative choice in TLS, where the redundant work is ephemeral,
negotiated, and invisible (Section~\ref{sec:hybrid-limits}). On chain the extra
signature is stored forever by everyone, cannot be negotiated away, and obliges
users to migrate twice. Importing the internet's default here doubles the
dominant cost.

\paragraph{Pitfall 5: quoting an exposure percentage without its cause.}
``One third of supply is exposed'' conflates outputs that publish a key by design
with addresses that were reused, and only the second is still growing and
reducible without a consensus change (Section~\ref{sec:two-exposures}).

\paragraph{Pitfall 6: treating migration as a scheduling problem.}
It is a throughput problem. Blockspace bounds it from below regardless of how
well the ecosystem coordinates (Section~\ref{sec:migration-throughput}), so a
plan that begins when Q-Day is announced cannot finish, however motivated
everyone is by then.

\section{Summary}
Blockchains sharpen the quantum threat and complicate the response:

\begin{itemize}
  \item the danger is to signatures (Shor), not hashing (Grover); an exposed public key on an immutable ledger is a permanent, ``store now, steal later'' liability --- which makes the shelf-life term in \textbf{Mosca's inequality} unbounded, so $x+y>z$ holds for any finite Q-Day and the only escape is to remove the exposure;
  \item exposure has \textbf{two causes} with different remedies: output types that publish a key by design, and address reuse --- the larger share, still accumulating, and reducible by wallet behaviour alone;
  \item \textbf{administrative keys} (stablecoin proxy admin, minter, pauser) are the concentrated case: few keys, shared blast radius, immune to multisig and hardware custody, but owned by someone who can actually act;
  \item \textbf{Bitcoin} proposes a new quantum-hardened output type (BIP-360) and a contested legacy-signature sunset (BIP-361), both draft soft forks;
  \item \textbf{Ethereum} pairs an emergency hard fork with STARK-based hash-preimage recovery and an account-abstraction migration path (EIP-7702 shipped, EIP-8141 draft), targeting first-layer post-quantum infrastructure in the 2028--2029 window;
  \item \textbf{Cosmos} has gone furthest in shipped code, registering ML-DSA-65 as an opt-in validator consensus key type, at $60\times$ the key size and $50\times$ the signature size of ed25519;
  \item a \textbf{layer 2} inherits all of these at once and can fix only some of them, so its plan must be organized by \emph{ownership} of each surface rather than by scheme;
  \item a symmetric or hash-based layer protects only what it \emph{carries}: Winternitz commitments, garbled circuits and FRI commitments have all been found wrapping elliptic-curve assertions, and a proof system can hide a discrete-log dependency in an argument that is not its headline commitment scheme;
  \item some exposures cannot be migrated at all --- an EVM pairing precompile has immutable callers, so it can be inventoried but never upgraded --- and aggregation has no post-quantum drop-in, which is why replacing BLS is an architectural project;
  \item \textbf{EdDSA chains} inherit a recovery route from RFC~8032's seed-then-hash key generation, letting an owner prove seed knowledge \emph{after} Q-Day with no prior action; ECDSA chains can imitate it per-wallet (BIP-32 already does) but cannot make it a consensus rule, because the structure is neither universal nor detectable on chain;
  \item migration is bounded from below by \textbf{blockspace}, not by willingness: at post-quantum witness sizes the floor is months to years of \emph{exclusive} blockspace, and years to decades at any share that leaves the chain usable --- a bound that follows from arithmetic rather than from any contested claim about quantum hardware;
  \item the timeline is uncertain but the required lead time is long, which is why these communities are debating the transition years before any quantum computer can execute the attack.
\end{itemize}

This is where cryptographic theory meets the messiest realities of deployment -- immutable history, decentralized governance, and irreversible loss -- and it is a fitting place to close: the mathematics is ready, and the harder work of migrating the systems that depend on it has only begun.

\printbibliography[heading=chapterrefs]
